%% toc: P2 17 17.1,17.2,17.3,17.7
\section{Proofs for first-order sequent adequacy}
\label{sec:annex-p02-w02b-p02-019-lk-nd}

Notation is that of \cref{sec:p02-w02b-p02-019-lk-nd}.  Two facts about the
syntax are used repeatedly.  First, \(\FV(\varphi)\subseteq
\operatorname{Var}(\varphi)\) for every formula: the clauses defining \(\operatorname{Var}\) in
\cref{sec:fol-syntax-semantics-03} are those defining \(\FV\), except that a
quantifier adds its bound variable instead of removing it, so the inclusion follows by induction
on formulas.  Second, \(\operatorname{Var}(\varphi)\) is finite and the
variable alphabet is infinite, so a variable outside any finite set of
variables exists; it is named, for definiteness, as the least such variable
in the fixed order \(x_0,x_1,\ldots\).

\begin{annexitem}{annex:p02-w02b-p02-019-renaming}{thm:p02-w02b-p02-019-renaming}{Renaming identities}
\begin{proof}
\emph{Terms.}  For a term \(s\) and variables \(x,u,w\) we show:
(a) \(s[x/x]=s\); (b) if \(u\notin\operatorname{Var}(s)\), then
\(\operatorname{Var}(s[u/x])\subseteq\operatorname{Var}(s)\cup\{u\}\) and
\((s[u/x])[w/u]=s[w/x]\); (c) if \(u\notin\operatorname{Var}(s)\), then
\(s[w/u]=s\).  Use induction on \(s\) with the clauses for substitution in a term in
\cref{sec:fol-syntax-semantics-04}.  If \(s\) is the variable \(x\), then
\(x[x/x]=x\); \(x[u/x]=u\), whose variable set is \(\{u\}\); and
\(u[w/u]=w=x[w/x]\).  If \(s\) is a variable \(v\ne x\), then
\(v[x/x]=v\) and \(v[u/x]=v\); since \(u\notin\operatorname{Var}(s)=\{v\}\),
\(v\ne u\), so \(v[w/u]=v=v[w/x]\).  Claim (c) at a variable \(v\) holds
because \(v\ne u\).  If \(s=f(s_1,\ldots,s_n)\) (including \(n=0\)), each
clause acts componentwise and
\(\operatorname{Var}(s)=\bigcup_i\operatorname{Var}(s_i)\), so all three
claims follow from the induction hypotheses for the \(s_i\).

\emph{Auxiliary claim.}  If \(u\notin\operatorname{Var}(\varphi)\) and
\(w\notin\operatorname{Var}(\varphi)\), then \(\varphi[w/u]=\varphi\).
Induct on the complexity of \(\varphi\).  For an atom the claim is (c) for
each argument term; \(\bot[w/u]=\bot\); a connective case follows from the
hypotheses for its immediate subformulas, whose variable sets are contained
in \(\operatorname{Var}(\varphi)\).  Let \(\varphi=Qv\,\chi\).  Since
\(v\in\operatorname{Var}(\varphi)\), \(v\ne u\) and \(v\ne w\), so
\(v\notin\operatorname{Var}(w)=\{w\}\).  The second clause of
\cref{def:p02-fol-capture-avoiding-substitution} applies and gives
\(Qv(\chi[w/u])\), which is \(Qv\,\chi\) by the induction hypothesis.

\emph{Item 1.}  Induct on the complexity of \(\varphi\).  Atoms follow from
(a), \(\bot\) is unchanged, and connectives act componentwise.  Let
\(\varphi=Qv\,\chi\).  If \(v=x\), the first substitution clause returns
\(Qx\,\chi\).  If \(v\ne x\), then \(v\notin\operatorname{Var}(x)=\{x\}\), the
second clause gives \(Qv(\chi[x/x])\), and the induction hypothesis gives
\(Qv\,\chi\).  The third clause never applies.

\emph{Item 2.}  Induct on the complexity of \(\varphi\), for all
\(x,u,w\) satisfying the hypotheses.  For an atom \(R(s_1,\ldots,s_n)\) or
\(s_1=s_2\), the hypotheses hold for each argument term, and (b) gives both
claims componentwise; the relation symbol, or the identity sign, is
unchanged.  For \(\bot\) both claims are trivial.  For \(\neg\chi\) and
\(\chi_1\circ\chi_2\), the hypotheses pass to the immediate subformulas,
because their variable sets are contained in that of \(\varphi\), and the
substitutions act componentwise.

Let \(\varphi=Qv\,\chi\).  Then \(v\in\operatorname{Var}(\varphi)\), so
\(u\ne v\) and \(w\ne v\); also \(u,w\notin\operatorname{Var}(\chi)\).

\emph{Case \(v=x\).}  The first substitution clause gives
\(\varphi[u/x]=Qx\,\chi=\varphi\), so the variable-set claim holds.  Next,
\((Qx\,\chi)[w/u]\): the binder \(x\) differs from \(u\) and is not in
\(\operatorname{Var}(w)=\{w\}\), so the second clause gives
\(Qx(\chi[w/u])\), which equals \(Qx\,\chi\) by the auxiliary claim.  On the
other side, \(\varphi[w/x]=(Qx\,\chi)[w/x]=Qx\,\chi\) by the first clause.

\emph{Case \(v\ne x\).}  Since \(v\ne u\), \(v\notin\operatorname{Var}(u)\),
so the second clause gives \(\varphi[u/x]=Qv(\chi[u/x])\).  Its variable set
is \(\{v\}\cup\operatorname{Var}(\chi[u/x])\), which by the induction
hypothesis is contained in
\(\{v\}\cup\operatorname{Var}(\chi)\cup\{u\}=\operatorname{Var}(\varphi)\cup\{u\}\).
Next, \((Qv(\chi[u/x]))[w/u]\): the binder \(v\) differs from \(u\) and from
\(w\), so the second clause applies again and gives
\(Qv((\chi[u/x])[w/u])\), which is \(Qv(\chi[w/x])\) by the induction
hypothesis.  Finally \(\varphi[w/x]=(Qv\,\chi)[w/x]\): here \(v\ne x\) and
\(v\ne w\), so the second clause gives \(Qv(\chi[w/x])\).  The two results are
literally the same formula.

\emph{Equality-freeness.}  We show that for every formula \(\varphi\), term
\(t\) and variable \(x\), an identity atom occurs in \(\varphi[t/x]\) exactly
when one occurs in \(\varphi\).  Induct on complexity, for all \(t,x\).  An
atom keeps its relation symbol or its identity sign; \(\bot\) and the
connectives are immediate.  At \(Qy\,\psi\), the first clause returns the
formula unchanged; the second gives \(Qy(\psi[t/x])\); the third gives
\(Qz((\psi[z/y])[t/x])\).  By \cref{lem:substitution-complexity},
\(\psi[z/y]\) has the complexity of \(\psi\), which is smaller than that of
\(Qy\,\psi\), so the induction hypothesis applies to \(\psi[z/y]\) and then
to \(\psi\).  The lemma's last sentence is the case in which \(t\) is a
variable.
\end{proof}
\end{annexitem}

\begin{annexitem}{annex:p02-w02b-p02-019-alpha-kernels}{thm:p02-w02b-p02-019-alpha-kernels}{Equality-free alpha derivations}
\begin{proof}
Let \(K\) be the set of pairs \((\varphi,\psi)\) of formulas with the
following two properties.
\begin{enumerate}
 \item[(K1)] \(\varphi\) is equality-free exactly when \(\psi\) is.
 \item[(K2)] If both are equality-free, \texttt{sheet:lk} derives
   \(\varphi\Rightarrow\psi\) and \(\psi\Rightarrow\varphi\) using only
   identity, weakening, the propositional and quantifier rules, and Cut,
   with every formula of the derivations equality-free.
\end{enumerate}
By \cref{def:p02-fol-alpha-equivalence}, \(=_{\alpha}\) is the least
equivalence relation that is a congruence for every formula constructor and
contains every generating pair.  We show that \(K\) has all of these
properties.  Then \({=_{\alpha}}\subseteq K\), which is the lemma.  In each
step below, the sequent rules are those of \texttt{sheet:lk}, applied with the
contexts named in parentheses; the formulas introduced are the given
formulas, their immediate subformulas, and their substitution instances by
variables, which are equality-free by
\cref{thm:p02-w02b-p02-019-renaming}.

\emph{Generating pairs.}  Let \(y\notin\operatorname{Var}(\varphi)\) and
consider \(Qx\,\varphi\) and \(Qy\,\chi\), where \(\chi\defeq\varphi[y/x]\).
(K1) holds by the last sentence of \cref{thm:p02-w02b-p02-019-renaming}.  For
(K2), choose a variable \(z\notin\operatorname{Var}(\varphi)\cup\{x,y\}\).
\Cref{thm:p02-w02b-p02-019-renaming}(2) with \(u\defeq y\) and \(w\defeq z\)
gives
\[
 \operatorname{Var}(\chi)\subseteq\operatorname{Var}(\varphi)\cup\{y\}
 \qquad\text{and}\qquad
 \chi[z/y]=\varphi[z/x].
\]
Write \(\theta\) for this common formula.  The variable \(z\) is not in
\(\operatorname{Var}(\varphi)\cup\{x\}\) or in
\(\operatorname{Var}(\chi)\cup\{y\}\), so \(z\) is free in neither
\(Qx\,\varphi\) nor \(Qy\,\chi\).  (When \(y=x\), the hypothesis says that
\(x\) does not occur in \(\varphi\); the same construction applies.)

For \(Q=\forall\), the derivation of \(\forall x\,\varphi\Rightarrow\forall
y\,\chi\) has three steps.
\begin{enumerate}
 \item Identity: \(\theta\Rightarrow\theta\).
 \item \(\forall L\) with principal formula \(\forall x\,\varphi\), term
   \(z\), and contexts \(\Gamma=\varnothing\), \(\Delta=\{\theta\}\).  Its
   premise is \(\varphi[z/x]\Rightarrow\theta\), which is step~1.  Its
   conclusion is \(\forall x\,\varphi\Rightarrow\theta\).
 \item \(\forall R\) with principal formula \(\forall y\,\chi\),
   eigenvariable \(z\), and contexts \(\Gamma=\{\forall x\,\varphi\}\),
   \(\Delta=\varnothing\).  Its premise is
   \(\forall x\,\varphi\Rightarrow\chi[z/y]\), which is step~2.  The side
   condition requires that \(z\) be free in no formula of
   \(\Gamma,\Delta,\forall y\,\chi\); this was checked above.  The
   conclusion is \(\forall x\,\varphi\Rightarrow\forall y\,\chi\).
\end{enumerate}
The reverse derivation is the same with the two quantified formulas
interchanged: identity \(\theta\Rightarrow\theta\); \(\forall L\) on
\(\forall y\,\chi\) with term \(z\), whose premise is
\(\chi[z/y]\Rightarrow\theta\); then \(\forall R\) on \(\forall x\,\varphi\)
with eigenvariable \(z\), whose premise is
\(\forall y\,\chi\Rightarrow\varphi[z/x]\).

For \(Q=\exists\), the derivation of \(\exists x\,\varphi\Rightarrow\exists
y\,\chi\) is:
\begin{enumerate}
 \item Identity: \(\theta\Rightarrow\theta\).
 \item \(\exists R\) with principal formula \(\exists y\,\chi\), term \(z\),
   and contexts \(\Gamma=\{\theta\}\), \(\Delta=\varnothing\).  Its premise
   is \(\theta\Rightarrow\chi[z/y]\), which is step~1.  Its conclusion is
   \(\theta\Rightarrow\exists y\,\chi\).
 \item \(\exists L\) with principal formula \(\exists x\,\varphi\),
   eigenvariable \(z\), and contexts \(\Gamma=\varnothing\),
   \(\Delta=\{\exists y\,\chi\}\).  Its premise is
   \(\varphi[z/x]\Rightarrow\exists y\,\chi\), which is step~2.  The side
   condition requires that \(z\) be free in no formula of
   \(\exists y\,\chi\) and \(\exists x\,\varphi\); this was checked above.
   The conclusion is \(\exists x\,\varphi\Rightarrow\exists y\,\chi\).
\end{enumerate}
The reverse derivation interchanges the two existential formulas in the same
way.  None of these derivations uses Cut.

\emph{Reflexivity.}  For every formula \(\varphi\), (K1) is trivial, and if
\(\varphi\) is equality-free, the identity sequent
\(\varphi\Rightarrow\varphi\) gives both derivations.

\emph{Symmetry.}  (K1) and (K2) are symmetric in \(\varphi\) and \(\psi\).

\emph{Transitivity.}  Let \((\varphi,\chi)\in K\) and \((\chi,\psi)\in K\).
(K1) for \((\varphi,\psi)\) follows by chaining the two equivalences.  If
\(\varphi\) and \(\psi\) are equality-free, so is \(\chi\), by (K1).  Cut
with cut formula \(\chi\) and contexts \(\Gamma=\{\varphi\}\),
\(\Delta=\Sigma=\varnothing\), \(\Pi=\{\psi\}\) takes the premises
\(\varphi\Rightarrow\chi\) and \(\chi\Rightarrow\psi\) to
\(\varphi\Rightarrow\psi\).  The reverse is the Cut of
\(\psi\Rightarrow\chi\) and \(\chi\Rightarrow\varphi\).  This is the only
step that introduces Cut.

\emph{Congruence.}  In each case (K1) holds because a compound formula is
equality-free exactly when its immediate subformulas are, and the
hypotheses supply (K1) for those.  For (K2), assume the compounds are
equality-free; then so are their immediate subformulas, and the hypotheses
supply their derivations.
\begin{enumerate}
 \item Negation.  From \(\psi\Rightarrow\varphi\), \(\neg L\) (contexts
   \(\Gamma=\{\psi\}\), \(\Delta=\varnothing\)) gives
   \(\psi,\neg\varphi\Rightarrow\), and \(\neg R\) (contexts
   \(\Gamma=\{\neg\varphi\}\), \(\Delta=\varnothing\)) gives
   \(\neg\varphi\Rightarrow\neg\psi\).  The reverse starts from
   \(\varphi\Rightarrow\psi\).
 \item Conjunction.  From \(\varphi_1\Rightarrow\psi_1\) and
   \(\varphi_2\Rightarrow\psi_2\), left weakening gives
   \(\varphi_1,\varphi_2\Rightarrow\psi_1\) and
   \(\varphi_1,\varphi_2\Rightarrow\psi_2\); \(\land R\) with the shared
   context \(\Gamma=\{\varphi_1,\varphi_2\}\) gives
   \(\varphi_1,\varphi_2\Rightarrow\psi_1\land\psi_2\); \(\land L\) gives
   \(\varphi_1\land\varphi_2\Rightarrow\psi_1\land\psi_2\).
 \item Disjunction.  From \(\varphi_1\Rightarrow\psi_1\), right weakening
   and \(\lor R\) give \(\varphi_1\Rightarrow\psi_1\lor\psi_2\); from
   \(\varphi_2\Rightarrow\psi_2\), the same steps give
   \(\varphi_2\Rightarrow\psi_1\lor\psi_2\), the succedent being a multiset.
   \(\lor L\) with the shared context \(\Delta=\{\psi_1\lor\psi_2\}\) gives
   \(\varphi_1\lor\varphi_2\Rightarrow\psi_1\lor\psi_2\).
 \item Implication.  From \(\psi_1\Rightarrow\varphi_1\), right weakening
   gives \(\psi_1\Rightarrow\psi_2,\varphi_1\).  From
   \(\varphi_2\Rightarrow\psi_2\), left weakening gives
   \(\psi_1,\varphi_2\Rightarrow\psi_2\).  These are the two premises of
   \(\to L\) with principal formula \(\varphi_1\to\varphi_2\) and contexts
   \(\Gamma=\{\psi_1\}\), \(\Delta=\{\psi_2\}\); its conclusion is
   \(\psi_1,\varphi_1\to\varphi_2\Rightarrow\psi_2\).  Then \(\to R\) gives
   \(\varphi_1\to\varphi_2\Rightarrow\psi_1\to\psi_2\).
 \item Quantifiers with the same bound variable.  By
   \cref{thm:p02-w02b-p02-019-renaming}(1), \(\varphi[x/x]=\varphi\) and
   \(\psi[x/x]=\psi\).  For \(\forall\): \(\forall L\) with term \(x\) takes
   \(\varphi\Rightarrow\psi\) to \(\forall x\,\varphi\Rightarrow\psi\); then
   \(\forall R\) with eigenvariable \(x\) and context
   \(\Gamma=\{\forall x\,\varphi\}\) gives
   \(\forall x\,\varphi\Rightarrow\forall x\,\psi\).  Its side condition
   holds because \(x\) is free in neither \(\forall x\,\varphi\) nor
   \(\forall x\,\psi\).  For \(\exists\): \(\exists R\) with term \(x\) takes
   \(\varphi\Rightarrow\psi\) to \(\varphi\Rightarrow\exists x\,\psi\); then
   \(\exists L\) with eigenvariable \(x\) and context
   \(\Delta=\{\exists x\,\psi\}\) gives
   \(\exists x\,\varphi\Rightarrow\exists x\,\psi\), with the same side
   condition check.
\end{enumerate}
Each converse uses the converse hypothesis derivations in the same way.  No
congruence step uses Cut.  Thus \(K\) is an equivalence relation, a
congruence for every constructor, and contains every generating pair, so it
contains \(=_{\alpha}\).
\end{proof}
\end{annexitem}

\begin{annexitem}{annex:p02-w02b-p02-019-maps}{thm:p02-w02b-p02-019-maps}{Two-way derivation maps}
\begin{proof}
\emph{Clause 2} is \cref{thm:p07-eqlk-nd}, with the same \(D_{\Delta}\), the
same support convention and the same two sheets.  Its proof,
Annex~\ref{annex:p7.eqlk.nd}, gives one construction for every rule of
\texttt{sheet:lk-eq}, including Mix and the two equality rules, and one for
every rule of \texttt{sheet:nd-fol}, including \(=I\), \(=E\) and
\(\alpha\); its alpha case uses the derivations of
\cref{thm:p07-eqlk-alpha}.

\emph{Clause 1.}  Compare \texttt{sheet:lk} with \texttt{sheet:lk-eq}: every
rule of the first (identity, \(\bot L\), the two weakenings and contractions,
the negation, conjunction, disjunction and implication rules, Cut, the four
quantifier rules and Mix) occurs on the second with the same premises,
conclusion and side condition.  The second has in addition only the
equality initial sequent \(\Rightarrow t=t\) and the rule \(=L\).

\emph{From sequents to natural deduction.}  Let \(\pi\) be a derivation of
\(\Gamma\Rightarrow\Delta\) on \texttt{sheet:lk}, all of whose formulas are
equality-free.  By the comparison, \(\pi\) is also a derivation on
\texttt{sheet:lk-eq}, and Annex~\ref{annex:p7.eqlk.nd} turns it into a
derivation of \(D_{\Delta}\) on \texttt{sheet:nd-fol} with open assumptions
in \(\operatorname{supp}(\Gamma)\).  That construction applies an equality
rule of natural deduction only in its cases for the equality initial sequent
(which becomes \(=I\)) and for \(=L\) (which uses \(=I\) and \(=E\) to derive
symmetry and then replacement).  Neither rule occurs in \(\pi\).  Its other
cases use assumptions, the propositional introduction and elimination rules,
RAA, \(\forall I\), \(\forall E\), \(\exists I\), \(\exists E\), \(\alpha\),
and the renaming of an already constructed derivation by a permutation of
the variables.  The renaming applies an equality rule only to an equality
instance already present in the derivation being renamed.  The formulas it
introduces are formulas of \(\pi\); finite disjunctions of them; \(\bot\);
the formulas \(\neg D\), \(\neg\neg A\) and \(\neg(D\lor A)\) built from
them; the conjunctions \(A\land C\) of the nonvacuous-discharge detour; and
images of all of these under permutations of the variables.  A permutation
of variables changes no relation symbol and creates no identity atom.  Hence
every formula of the result is equality-free, and the result is a derivation
in the equality-free restriction of \texttt{sheet:nd-fol}.  This includes
the Mix rule and the freshening at \(\forall R\) and \(\exists L\), which
Annex~\ref{annex:p7.eqlk.nd} treats explicitly.

\emph{From natural deduction to sequents.}  Let \(\nu\) be a derivation of
\(D_{\Delta}\) in the equality-free restriction of \texttt{sheet:nd-fol},
with open assumptions in \(\operatorname{supp}(\Gamma)\).  Follow the
reverse construction of Annex~\ref{annex:p7.eqlk.nd}, which assigns to each
node with open-assumption support \(S\) and conclusion \(\chi\) a sequent
derivation of \(S^{\circ}\Rightarrow\chi\); here \(S^{\circ}\) has one
occurrence of each member of \(S\), and temporary discharge formulas are
added as that construction prescribes.  Its rule cases use identity,
weakening, contraction, \(\bot L\), the left and right propositional rules,
\(\forall L\), \(\forall R\), \(\exists L\), \(\exists R\) and Cut, all
present on \texttt{sheet:lk}.  At \(\forall I\) the eigenvariable of
\(\forall R\) is the generalized variable \(x\) itself; this is correct
because \(\varphi[x/x]=\varphi\) by
\cref{thm:p02-w02b-p02-019-renaming}(1).  The equality initial sequent and
the replacement derivation are used only in the cases \(=I\) and \(=E\),
which do not occur in \(\nu\).  The one remaining case is \(\alpha\), and we
write it out in full.

Let the node be an application of \(\alpha\) with premise
\(\Gamma_0\vdash\varphi\), conclusion \(\Gamma_0\vdash\psi\) and
\(\varphi=_{\alpha}\psi\), where all formulas are equality-free.  The rule
discharges nothing, so the open-assumption support \(S\) of the conclusion is
that of the premise.
\begin{enumerate}
 \item The construction applied to the premise subderivation gives a
   derivation on \texttt{sheet:lk} of \(S^{\circ}\Rightarrow\varphi\).
 \item \Cref{thm:p02-w02b-p02-019-alpha-kernels} gives a derivation on
   \texttt{sheet:lk} of \(\varphi\Rightarrow\psi\), all of whose formulas are
   equality-free.
 \item Cut with cut formula \(\varphi\) and contexts \(\Gamma=S^{\circ}\),
   \(\Delta=\varnothing\), \(\Sigma=\varnothing\), \(\Pi=\{\psi\}\) has
   premises \(S^{\circ}\Rightarrow\varphi\) (step~1) and
   \(\varphi\Rightarrow\psi\) (step~2), and conclusion
   \(S^{\circ}\Rightarrow\psi\).
\end{enumerate}
Cut has no side condition, and no eigenvariable condition is involved.  The
conclusion \(S^{\circ}\Rightarrow\psi\) is exactly the sequent the
construction requires at this node.  This replaces the use of
\cref{thm:p07-eqlk-alpha}, whose derivations are on \texttt{sheet:lk-eq}, by
derivations on \texttt{sheet:lk}.

At the root the construction gives
\(\operatorname{supp}(\Gamma)^{\circ}\Rightarrow D_{\Delta}\).  The
derivation of \(D_{\Delta}\Rightarrow\Delta\) in
Annex~\ref{annex:p7.eqlk.nd} uses identity at each disjunct, \(\lor L\) at
each internal node of the disjunction, \(\bot L\) when
\(\operatorname{supp}(\Delta)\) is empty, and right weakening for repeated
succedent formulas; a Cut and left weakening then give
\(\Gamma\Rightarrow\Delta\).  All of these rules are on \texttt{sheet:lk},
and every formula is a formula of \(\nu\), of \(\Gamma\) or \(\Delta\), or of
a derivation from step~2, so every formula is equality-free.

Both directions of clause~1 are proved.  The constructions are recursions on
finite trees and make no recursive call on a derivation they have built.
\end{proof}
\end{annexitem}

\begin{annexitem}{annex:p02-w02b-p02-019-choice-free}{thm:p02-w02b-p02-019-choice-free}{Choice-free LK adequacy}
\begin{proof}
\emph{Soundness.}  A Mix step with premises
\(\Gamma\Rightarrow\Delta,\varphi^{m}\) and
\(\varphi^{n},\Sigma\Rightarrow\Pi\), \(m,n\ge 1\), is replaced by \(m-1\)
right contractions, \(n-1\) left contractions and one Cut on \(\varphi\);
the conclusion \(\Gamma,\Sigma\Rightarrow\Delta,\Pi\) is the same.  After
this replacement every rule of a \texttt{sheet:lk} derivation is a rule of
the calculus of \cref{thm:fol-lk-soundness}, so every derivable sequent is
satisfied by every structure and assignment, that is, valid.  No choice
principle is used.

\emph{Completeness.}  Let \(\Gamma\Rightarrow\Delta\) be valid over all
\(\Lang\)-structures.  Let \(\Lang_0\) consist of the function symbols
(including constants) and relation symbols that occur in the sequent.  The
sequent is a finite string, so \(\Lang_0\) is finite, possibly empty, and
every formula of the sequent is an \(\Lang_0\)-formula.

First, the sequent is valid over all \(\Lang_0\)-structures.  Let
\(\mathcal M_0\) be an \(\Lang_0\)-structure and \(s\) an assignment in it.  Its
domain is nonempty; fix one element \(a\).  Define an \(\Lang\)-structure
\(\mathcal M\) with the same domain: each symbol of \(\Lang_0\) has its
interpretation in \(\mathcal M_0\); each function symbol of \(\Lang\) outside
\(\Lang_0\) of arity \(n\) is interpreted as the function with constant value
\(a\) (for \(n=0\), the constant \(a\)); each relation symbol outside
\(\Lang_0\) is interpreted as the empty relation (for arity \(0\), as false).
This single explicit definition uses one fixed element and no choice from a
family of sets.  Then \(\mathcal M\mathbin{\upharpoonright}\Lang_0=\mathcal M_0\), and \(s\) is an
assignment in \(\mathcal M\).  By validity, \(\mathcal M,s\) fails some member of
\(\operatorname{supp}(\Gamma)\) or satisfies some member of
\(\operatorname{supp}(\Delta)\).  By \cref{lem:reduct-invariance}, applied to
the extension \(\Lang\) of \(\Lang_0\), each formula of the sequent has the
same truth value at \(\mathcal M,s\) and at \(\mathcal M_0,s\).  So \(\mathcal M_0,s\)
satisfies the sequent.

Second, \(\Lang_0\) is finite, so the finite-or-countable clause of
\cref{thm:p07-canonical-lk-completeness}, which needs no Choice, gives a
derivation of \(\Gamma\Rightarrow\Delta\) on \texttt{sheet:lk} over
\(\Lang_0\).

Third, that derivation is a derivation over \(\Lang\).  Its formulas are
\(\Lang\)-formulas, and the terms used in \(\forall L\) and \(\exists R\) are
\(\Lang\)-terms.  The substitution operation, \(\operatorname{Var}\) and
\(\FV\) are defined by recursion on the formula and on the fixed variable
alphabet, not on the signature, so each substitution instance and each
eigenvariable condition reads the same over \(\Lang\).

Neither step uses a choice principle, and the derivation obtained may
contain Cut, because the cited theorem's proof uses Cut.
\end{proof}
\end{annexitem}
